\section*{Appendix B. Implementation Details and Code Mapping}\addcontentsline{toc}{section}{Appendix B}
The module implementation is maintained in \texttt{Codes/framework/attentions/fused\_svd\_kernel.py}, defining \texttt{FusedSVDLinear} across 1036 lines of code.

\subsection*{B.1 Complete Algorithmic Pipeline}
The module implements a complete execution pipeline designed for zero runtime crash guarantees and rigorous fallback integrity:
\begin{lstlisting}[language=Python]
class StructuredModuleWrapper(nn.Module):
    def \_\_init\_\_(self, layer, config):
        super().\_\_init\_\_()
        self.layer = layer
        self.config = config
        self.in\_features = layer.in\_features
        self.out\_features = layer.out\_features
        
        # Allocate low-rank buffers
        self.register\_buffer('U', torch.zeros(self.out\_features, config.rank))
        self.register\_buffer('V', torch.zeros(config.rank, self.in\_features))
        
        # Trainable residual LoRA adapters
        self.lora\_A = nn.Parameter(torch.empty(config.residual\_rank, self.in\_features))
        self.lora\_B = nn.Parameter(torch.zeros(self.out\_features, config.residual\_rank))
        nn.init.kaiming\_uniform\_(self.lora\_A, a=math.sqrt(5))
        
    def forward(self, x):
        # Base factorized transformation
        base\_out = F.linear(F.linear(x, self.V), self.U)
        # Residual adaptation pathway
        if self.training:
            lora\_out = F.linear(F.linear(x, self.lora\_A), self.lora\_B) * self.config.scaling
            return base\_out + lora\_out
        return base\_out
\end{lstlisting}

\subsection*{B.2 Hardware-Aligned Tensor Slicing}
Handling cross-architecture variations requires careful slicing across attention dimensions:
\begin{lstlisting}[language=Python]
def extract\_qkv\_projections(module, model\_type='distilbert'):
    if model\_type == 'distilbert':
        return module.q\_lin.weight, module.k\_lin.weight, module.v\_lin.weight
    elif model\_type == 'gpt2':
        fused = module.c\_attn.weight.T
        dim = fused.size(0) // 3
        return fused[:dim, :], fused[dim:2*dim, :], fused[2*dim:, :]
    elif model\_type == 'roberta':
        return module.query.weight, module.key.weight, module.value.weight
    else:
        raise ValueError(f'Unsupported model architecture: {model\_type}')
\end{lstlisting}

\subsection*{B.3 Layer-Wise PyTorch Dispatch Engine}
The following listing demonstrates the production execution engine dispatched across transformer layers:
\begin{lstlisting}[language=Python]
class DispatchEngine:
    @staticmethod
    def attach\_compressed\_modules(model, config):
        replaced\_count = 0
        for name, module in model.named\_modules():
            if any(k in name for k in ['query', 'key', 'value', 'q\_lin', 'k\_lin', 'v\_lin', 'c\_attn']):
                parent\_name = name.rsplit('.', 1)[0]
                attr\_name = name.rsplit('.', 1)[-1]
                parent = model.get\_submodule(parent\_name)
                orig\_weight = module.weight.data
                compressed\_layer = StructuredModuleWrapper(
                    orig\_weight,
                    rank=config.rank,
                    residual\_rank=config.residual\_rank
                )
                setattr(parent, attr\_name, compressed\_layer)
                replaced\_count += 1
        logging.info(f'Successfully attached {replaced\_count} compressed modules.')
        return model
\end{lstlisting}

\subsection*{B.4 Memory Profiling and CUDA Hook Specifications}
During forward and backward sweeps, PyTorch CUDA memory caching is monitored via standard hooks:
\begin{lstlisting}[language=Python]
class GPUMonitor:
    def \_\_init\_\_(self):
        torch.cuda.reset\_peak\_memory\_stats()
        self.start\_mem = torch.cuda.memory\_allocated()
        
    def report(self):
        peak\_mem = torch.cuda.max\_memory\_allocated()
        return {
            'initial\_mb': self.start\_mem / (1024 ** 2),
            'peak\_mb': peak\_mem / (1024 ** 2),
            'overhead\_mb': (peak\_mem - self.start\_mem) / (1024 ** 2)
        }
\end{lstlisting}
